\documentclass{article}

\usepackage{amsmath, amsthm, amssymb, amsfonts}
\usepackage{thmtools}
\usepackage{graphicx}
\usepackage{setspace}
\usepackage{geometry}
\usepackage{float}
\usepackage{hyperref}
\usepackage[utf8]{inputenc}
\usepackage[english]{babel}
\usepackage{framed}
\usepackage[dvipsnames]{xcolor}
\usepackage{tcolorbox}

\colorlet{LightGray}{White!90!Periwinkle}
\colorlet{LightOrange}{Orange!15}
\colorlet{LightGreen}{Green!15}

\newcommand{\HRule}[1]{\rule{\linewidth}{#1}}

\declaretheoremstyle[name=Theorem,]{thmsty}
\declaretheorem[style=thmsty,numberwithin=section]{theorem}
\tcolorboxenvironment{theorem}{colback=LightGray}

\declaretheoremstyle[name=Proposition,]{prosty}
\declaretheorem[style=prosty,numberlike=theorem]{proposition}
\tcolorboxenvironment{proposition}{colback=LightOrange}

\declaretheoremstyle[name=Principle,]{prcpsty}
\declaretheorem[style=prcpsty,numberlike=theorem]{principle}
\tcolorboxenvironment{principle}{colback=LightGreen}

\setstretch{1.2}
\geometry{
    textheight=9in,
    textwidth=5.5in,
    top=1in,
    headheight=12pt,
    headsep=25pt,
    footskip=30pt
}

% ------------------------------------------------------------------------------

\begin{document}

% ------------------------------------------------------------------------------
% Cover Page and ToC
% ------------------------------------------------------------------------------

\title{ \normalsize \textsc{}
		\\ [2.0cm]
		\HRule{1.5pt} \\
		\LARGE \textbf{\uppercase{Relatório do trabalho\\Padaria}
		\HRule{2.0pt} \\ [0.6cm] \LARGE{Ferramentas da Internet} \vspace*{10\baselineskip}}
		}
\date{}
\author{\textbf{Autores:}
		Gabrieli Pio\\
		Amanda\\
        Beatriz\\
        Gabriel Arthur\\
        Gabriel Vitor\\
        Daniel\\
		Mariana, 2026}

\maketitle
\newpage

% ------------------------------------------------------------------------------

\section{Introdução}
O minimundo é a descrição textual e simplificada de uma parte da realidade que se deseja transformar em um banco de dados. Ele funciona como um recorte do "mundo real", detalhando as regras, os objetos e os processos específicos de um negócio para que o analista saiba exatamente quais informações devem ser capturadas e como elas se relacionam no sistema.
\begin{itemize}
   
\end{itemize}

\section{Minimundo}
Uma padaria necessita de um sistema de banco de dados para sua gerência com informações dos cardápios, produtos, ingredientes, pedidos, clientes, funcionários, fornecedores, compras, vendas e gerente. A padaria possui cardápios, identificados por um código, contendo informações de data, dia, hora e duração. Cada cardápio inclui diversos produtos, sendo que um mesmo produto pode estar presente em diferentes cardápios. Os produtos são identificados por um número de prato e possuem nome, descrição e preço. Cada produto pode ser composto por vários ingredientes, e um ingrediente pode compor diferentes produtos. Os ingredientes possuem código, nome, descrição e validade. Eles são adquiridos por meio de compras. Cada compra é identificada por um código e possui data, dia, hora e valor. Uma compra pode incluir vários ingredientes, e cada ingrediente pode estar presente em várias compras. As compras são realizadas junto a fornecedores, que possuem identificação, nome, endereço e telefone.Um fornecedor pode estar associado a várias compras, enquanto cada compra está vinculada a um único fornecedor. Os clientes são identificados por CPF e possuem nome, telefone e e-mail. Um cliente pode realizar vários pedidos, mas cada pedido está associado a apenas um cliente. Os pedidos possuem identificação, número, data, dia, hora e valor. Cada pedido pode conter vários produtos, e um produto pode estar presente em diversos pedidos. Os pedidos são atendidos por funcionários, que possuem identificação, nome, cargo, salário, telefone e e-mail. Um funcionário pode atender vários pedidos, mas um pedido é atendido por um único funcionário.Os pedidos fazem parte de vendas, que possuem identificação do pedido, valor, data, dia e hora. Cada venda está associada a um pedido, e os funcionários podem listar diversas vendas. O banco de dados também contempla o gerente, que possui CPF, nome, telefone e e-mail. O gerente é responsável por monitorar os pedidos e as vendas realizadas no restaurante, além de liderar os outros funcionários.

\section{Conclusão}
Em resumo, esse mini mundo descreve um sistema de padaria.
O objetivo final não é só armazenar dados, mas saber gerência de forma prática um negócio . É a tecnologia garantindo que o sistema organize tudo de uma forma correta\newpage


\end{document}
